\begin{textsample}{POS Dim 2 – mg – Score 26.00 – mg/mg\_em\_2.txt}  \label{ex:f2_pos_195}
O processo de construção do \textbf{Currículo} Referência de Minas Gerais Em um momento histórico de definição de uma Base Nacional Comum Curricular e o desafio da elaboração de um \textbf{Currículo} Referência que \textbf{atenda} a todo o estado, em razão da sua extensão e número de \textbf{municípios} e escolas, o \textbf{Regime} de Colaboração passa a ser central no cenário educacional.
Em um esforço conjunto para reunir a imensa “ Minas Gerais ” e construir um documento coletivo, a seccional da União Nacional dos \textbf{Dirigentes} \textbf{Municipais} de Educação – UNDIME/MG, as escolas \textbf{privadas} de educação e a Secretaria de Estado de Educação – SEE-MG colaboraram lado a lado para a redação desse documento.
Considerando, fundamentalmente, que os estudantes transitam entre as \textbf{redes} ao longo de seu percurso escolar, ora em escolas \textbf{municipais}, ora em \textbf{estaduais}, ora em instituições privadas, bem como transitam entre os territórios, reforça -se aí a importância de uma parte comum nos \textbf{currículos}.
Se o documento curricular pretende garantir, de fato, os direitos de aprendizagem, torna -se impossível fragmentar a vida escolar de nossos estudantes e, assim, em colaboração, buscamos garantir \textbf{trajetórias} de sucesso acadêmico, somadas ao desenvolvimento integral das crianças, \textbf{adolescentes}, jovens e adultos que estão e estarão na educação pública ou particular, buscando, por meio de instâncias decisórias estabelecidas, o diálogo permanente entre União, Estado e Municípios.
Minas Gerais é o estado brasileiro com maior número de \textbf{municípios} ( 853 ), representando 15 \% do total do país ( 5570 \textbf{municípios} ).
Apresenta um retrato quase sempre fiel da realidade brasileira, com 10 \% ( 20.7 milhões ) da população nacional ( 209.3 milhões ), representando a grande diversidade regional, econômica, política e social.
Em termos educacionais, nosso estado conta com 16.151 escolas, sendo 3.622 \textbf{estaduais}, 8.751 \textbf{municipais} e 3.778 privadas, distribuídas em 47 regionais de ensino ( SRE ), e com 4.032.949 estudantes matriculados, sendo que 86 \% deles estão na \textbf{rede} pública.
Com a maioria das escolas e das matrículas pertencentes à \textbf{rede} pública, garantir uma educação de \textbf{qualidade} com \textbf{equidade} é princípio norteador das \textbf{políticas} públicas de educação nas \textbf{redes} \textbf{municipal}, \textbf{estadual} e particular.
Localizado na região Sudeste, o estado de Minas Gerais é o quarto em área territorial do país com 586.520,732 km²¹.
Com uma população estimada de 21.119.536 habitantes², distribuídos em 853 \textbf{municípios}, é o segundo mais populoso, o que contribui significativamente para sua grande pluralidade regional.
Sua diversidade regional é resultado de um processo histórico de ocupação do território marcado por diferentes fatores, desde aqueles de ordem socioeconômica até os naturais de clima e vegetação.
Essa diversidade se traduz no que podemos entender como várias “ Minas Gerais ” dentro dos limites do estado, exigindo, portanto, diferentes formas de abordagem e atuação sobre a realidade mineira.
De fato, a efetividade de qualquer iniciativa parte necessariamente da compreensão dessa realidade.
A figura abaixo mostra a atual divisão do Estado de Minas Gerais em 17 Territórios de Desenvolvimento, marcados pelas suas semelhanças regionais.
As semelhanças que agrupam os \textbf{municípios} conduzem para a formação de \textbf{perfis} razoavelmente consistentes, para cada território, com características aferíveis e, principalmente, comparáveis.
Nesse sentido, é possível vislumbrar as diferenças existentes, e consequentemente, perceber a dificuldade do estado em \textbf{atender} o Ensino Médio por meio de \textbf{políticas} públicas que considerem essa diversidade regional, e que garantam efetivamente o acesso e a permanência das juventudes a essa etapa, bem como o sucesso na conclusão da Educação Básica.
Essa diversidade no desenvolvimento dos territórios mineiros é analisada considerando elementos como : os diferentes graus de urbanização, dado pela proporção da população urbana na população total ; o Índice de Desenvolvimento Humano Municipal ( IDH-M ), que utiliza a mesma lógica do IDH Global, mas com uma metodologia mais adequada à realidade dos \textbf{municípios}, sustentando -se em três dimensões – a renda, a educação e a longevidade.
E também, elementos como a distribuição das turmas entre as zonas rural e urbana ; a formação dos docentes ; a diversidade no atendimento promovido pela \textbf{rede} pública.
E, finalmente, para além dos regionalismos, Minas Gerais apresenta uma importante demanda por atendimentos especializados, traduzida nas instituições que se encontram em localizações diferenciadas : as áreas de assentamento, de uso sustentável, quilombolas e terras indígenas.
Portanto, a significativa heterogeneidade das demandas de Minas Gerais aponta para a necessidade de se considerar as diferenças sociais, econômicas e demográficas apresentadas, no planejamento e na execução de \textbf{políticas} públicas, garantindo que a \textbf{oferta} da Educação Básica seja ampliada e se efetive com a \textbf{qualidade} social desejada, respeitando as particularidades de cada território.
Para a elaboração do novo \textbf{currículo} de Minas Gerais que \textbf{atenda} tamanha diversidade foi estabelecido um modelo de governança dinâmico, capaz de lidar com as particularidades do estado e com as diversas \textbf{entidades} que atuam diretamente para melhoria da educação pública e particular.
Foi estabelecida uma Comissão Estadual, com representações políticas de \textbf{órgãos} e \textbf{entidades}, um Comitê Executivo para condução e tomada de decisão, uma Coordenação \textbf{Técnica} para encaminhamento dos trabalhos e Grupos de Trabalho de \textbf{Currículo} para redação do documento, conforme Figura 2.
Na elaboração deste documento, o \textbf{regime} de colaboração efetivou -se na formação dos Grupos de Trabalho de \textbf{Currículo} e, sobretudo, nos inúmeros momentos de discussão em que \textbf{profissionais} de diversas áreas do conhecimento, oriundos das várias regiões do estado, reuniram -se para discutir o \textbf{Currículo} Referência de Minas Gerais, de modo a conferir -lhe um caráter próprio, incorporando as diretrizes e normativas da BNCC, bem como dos \textbf{preceitos} de uma educação libertadora, que vise à \textbf{equidade} e à \textbf{qualidade} educacional dos sistemas de ensino, promovendo a inclusão, reconhecendo e valorizando as diversidades.
O primeiro movimento de aproximação das escolas de Minas Gerais com a BNCC \textbf{homologada} foi com a realização do Dia D, um momento de reflexão e engajamento de toda a Rede Pública, visando a garantir aos \textbf{profissionais} a oportunidade de contribuir para as definições conceituais do documento curricular \textbf{preliminar}.
Os Grupos de Trabalho de \textbf{Currículo} ( Educação Infantil, Ensino Fundamental e Médio ), responsáveis pela entrega do documento \textbf{preliminar}, participaram de encontros \textbf{formativos} realizados pelo Ministério da Educação – MEC e por formações internas oferecidas por professores e pesquisadores de universidades mineiras.
Concomitantemente, essa equipe realizou estudos aprofundados da Base Nacional Comum Curricular e de diversos documentos curriculares de Minas Gerais e de outros estados que possibilitaram definições conceituais e a elaboração do formato do documento atual.
Minas Gerais avança ao propor um \textbf{currículo} de referência que coloca as crianças, as juventudes e adultos no centro do processo de ensino e de aprendizagem, considerando os sujeitos numa visão integral ; que dialoga com anseios diversos e múltiplas necessidades de formação ; que reverbera o processo de aprendizagem de forma participativa e produtora de conhecimentos, imanente às realidades dos sujeitos participantes ; que inova numa visão de formação para além dos conteúdos escolares e, também, para as práticas nas relações sociais no e com o mundo.
Nessa medida, o \textbf{Currículo} Referência é fruto do trabalho coletivo de centenas de \textbf{profissionais} de várias regiões do estado, versando sobre pluralidade de ideias, identidades e expressões mineiras e, em \textbf{conformidade} com a BNCC, sendo a referência curricular para as escolas de Educação Infantil e Ensino Fundamental desde 2019 e para o Ensino Médio de todo o estado a partir de 2021.

% matched lemmas: adolescente, atender, conformidade, currículo, dirigente, entidade, equidade, estadual, formativo, homologar, municipal, município, oferta, perfil, política, preceito, preliminar, privar, profissional, qualidade, rede, regime, trajetória, técnico, órgão
\end{textsample}
